\documentclass[12pt,a4paper]{article}
\usepackage[a4paper,margin=18mm]{geometry}
\usepackage[T1]{fontenc}
\usepackage{amsmath}
\usepackage{amssymb}
\usepackage{enumitem}

\pagestyle{empty}
\setlength{\parindent}{0pt}
\setlength{\parskip}{4pt}

\begin{document}

\begin{center}
  {\LARGE\bfseries Constructions and Loci}\\[2pt]
  {\Large Answer Key}
\end{center}

\begin{enumerate}[label=\textbf{Q\arabic*.}, leftmargin=*, itemsep=4pt, topsep=4pt]

\item Perpendicular bisector of \(AB\), passing through the midpoint of \(AB\). With \(AB=8\,\text{cm}\), it crosses \(AB\) \(4\,\text{cm}\) from \(A\) and \(B\).

\item \(3\,\text{cm}\).

\item The line through \(Q\) perpendicular to the given line.

\item \(3\,\text{cm}\).

\item \(M\) is the midpoint of \(CD\), and \(CM=3.5\,\text{cm}\).

\item The horizontal line through \(S\), perpendicular to the given vertical line.

\item The two constructed perpendiculars are parallel.

\item (a) Foot of perpendicular is \(2\,\text{cm}\) from \(A\); (b) height \(=4\,\text{cm}\); (c) area \(=14\,\text{cm}^2\).

\item (c) \(EG=\sqrt{4^2+4^2}=4\sqrt2\approx5.66\,\text{cm}\), measured as about \(5.7\,\text{cm}\).

\item \(O\) is the intersection of the perpendicular bisectors of the chords; radius \(=3.5\,\text{cm}\).

\item (b) A circle; (c) draw the circle centre \(C\), radius \(3\,\text{cm}\); (d) shade the inside of this circle.

\item (c) The midpoint of \(AB\); (d) the perpendicular bisector of \(AB\); (e) shade the half-plane on \(A\)'s side of the perpendicular bisector.

\item (a) \(3.5\,\text{cm}\); (b) circle centre \(L\), radius \(3.5\,\text{cm}\); (c) perpendicular bisector of \(LM\); (d) shade inside the circle and on \(M\)'s side of the perpendicular bisector.

\item (b) \(3.5\,\text{cm}\); (c) perpendicular; (d) \(3.5\times100=350\,\text{m}\).

\item The locus is the pair of parallel lines \(2\,\text{cm}\) above and \(2\,\text{cm}\) below the given line.

\item The locus is two line segments parallel to \(XY\) at distance \(2\,\text{cm}\), joined by semicircles of radius \(2\,\text{cm}\) centred at \(X\) and \(Y\). Priya is wrong because the ends are semicircular, not straight sides of a rectangle.

\item Shade the region that is more than \(3\,\text{cm}\) from the house wall, inside the circle centre \(P\) radius \(4\,\text{cm}\), and on the \(T\)-side of the perpendicular bisector of \(TB\). Approximately, along the \(3\,\text{cm}\) line from the wall the region runs from \(2.7\,\text{cm}\) to \(5.3\,\text{cm}\) from the left-hand fence; along the top fence it runs from \(3.6\,\text{cm}\) to \(7.5\,\text{cm}\) from the left-hand corner.

\item (a) Perpendicular bisector of \(AB\); (b) perpendicular bisector of \(BC\); (c) \(M\) is their intersection; (d) the perpendicular bisector of \(AC\) also passes through \(M\); (e) because \(MA=MB\) and \(MB=MC\), hence \(MA=MC\); (f) \(4\,\text{cm}\) on the map represents \(8\,\text{km}\) in real life.

\item (a) Intersections of the perpendicular bisector of \(OW\) with the circle centre \(S\) radius \(4\,\text{cm}\): approximately \(T_1=(6.2,3.9)\) and \(T_2=(11.1,4.4)\) cm from the bottom-left corner. (b) For example: the treasure is less than \(50\,\text{m}\) from \(O\), i.e. inside the circle centre \(O\) radius \(5\,\text{cm}\); this selects \(T_1\). (c) For example: circles of radii \(8.0\,\text{cm}\) about \(O\), \(6.5\,\text{cm}\) about \(W\), and \(5.3\,\text{cm}\) about \(S\). The first two meet at \(X\) and at a second point off the left edge; the third selects \(X\). If the map edge is allowed, the first two clues already select \(X\) alone on the map.

\end{enumerate}

\end{document}